\section{Nef and big canonical bundle: parallelism and integrability}
\label{sec:canonical}

Throughout this section, \(X\) and \(T_X=E_1\oplus E_2\) satisfy the
hypotheses of Corollary~\ref{thm:nef-big}; integrability is not assumed.
We prove it by an estimate for the second fundamental forms of the
summands. This estimate also gives parallelism on the regular canonical
locus, the additional input needed for the alternative proof in the
appendices.

\begin{proposition}\label{canon:contraction}
There is a projective birational morphism with connected fibers
\[
 f:X\longrightarrow Z
\]
onto a normal projective variety with canonical singularities and ample
\(\mathbf Q\)-Cartier canonical divisor.  Moreover,
\(K_X=f^*K_Z\), and \(f\) is an isomorphism over \(Z_{\mathrm{reg}}\).
\end{proposition}

\begin{proof}
The base-point-free theorem applies to \(D=K_X\): \(D\) is nef and
\(aD-K_X=(a-1)K_X\) is nef and big for an integer \(a>1\)
\cite[Theorem~3.3]{KM98}.  The associated contraction, with its Stein
factorization, has normal projective target and connected fibers.
For some positive integer \(m\) and ample line bundle \(L\) on \(Z\),
\[
 \mathcal O_X(mK_X)\simeq f^*L.
\]
Bigness makes \(f\) generically finite; connected fibers then make it
birational.  A proper birational morphism onto a normal variety is an
isomorphism over an open set whose complement has codimension at least
two.  On this open set \(L\) agrees with
\(\mathcal O_Z(mK_Z)\).  Reflexive extension therefore identifies
\(L\) with \(\mathcal O_Z(mK_Z)\) everywhere.  Thus \(K_Z\) is
\(\mathbf Q\)-Cartier and ample.

Choose compatible canonical divisors.  The divisor
\(K_X-f^*K_Z\) is exceptional and \(\mathbf Q\)-linearly trivial.
A principal exceptional divisor is zero: the corresponding rational
function has no divisor on the normal projective variety \(Z\), hence
is constant.  Consequently \(K_X=f^*K_Z\).  The discrepancies on any
higher resolution are nonnegative because \(X\) is smooth, so \(Z\)
has canonical, and in particular klt, singularities.

Over \(Z_{\mathrm{reg}}\), the relative Jacobian has zero divisor.
It is therefore nowhere zero, and \(f\) is a local biholomorphism
there.  Properness and birationality imply that this restriction is an
isomorphism.
\end{proof}

Set
\[
 X^0=f^{-1}(Z_{\mathrm{reg}}).
\]
This is a connected dense open subset of \(X\); we identify it with
\(Z_{\mathrm{reg}}\).  Our convention is
\(\Ric(\omega)=-\sqrt{-1}\,\partial\bar\partial\log\omega^n\),
so \([\Ric(\omega)]=2\pi c_1(X)\).

\begin{lemma}[The canonical metric]\label{canon:metric-input}
Fix a Kähler form \(\omega_A\) from a polarization of \(X\).
There are smooth Kähler metrics \(g_\varepsilon\) on \(X\), for
\(0<\varepsilon<1\), whose forms satisfy
\[
 \Ric(\omega_\varepsilon)
       =-\omega_\varepsilon+\varepsilon\omega_A,
 \qquad
 [\omega_\varepsilon]
       =2\pi c_1(K_X)+\varepsilon[\omega_A].
\]
As \(\varepsilon\to0\), they converge smoothly on compact subsets of
\(X^0\) to a Kähler metric \(g\) with \(\Ric(g)=-g\).
The completion of the intrinsic Riemannian distance of
\((X^0,g)\) is naturally homeomorphic to \(Z\).
\end{lemma}

\begin{proof}
Proposition~\ref{canon:contraction} makes \(f\) a smooth projective
crepant resolution of a canonical model.  The completion statement is
\cite[Theorem~1.2]{Song14}.  The approximating metrics and their smooth
convergence are \cite[Equation~(4.28) and Lemma~4.1]{Song14}, with
\(\varepsilon=e^{-t}\).  The cohomology identity follows directly from
the Ricci equation.
\end{proof}

\begin{proposition}\label{canon:parallel}
The connection of \(g\) preserves both \(E_1|_{X^0}\) and
\(E_2|_{X^0}\).  Both subbundles are integrable on all of \(X\).
They admit simultaneous local holomorphic product coordinates.
\end{proposition}

\begin{proof}
We use the subbundle-curvature method of approximate
K\"ahler--Einstein geometry, developed in the stability results of
Enoki \cite{Enoki88} and used for parallel splittings by Guenancia
\cite[Section~3.4, Step~4, Equations~(3.12)--(3.13)]{Guenancia16}.
Here the equality of the two summands' total first Chern class with
$c_1(T_X)$ gives a direct cancellation.
Give \(T_X\) the Hermitian metric induced by \(g_\varepsilon\).
Let \(P_{j,\varepsilon}\) be orthogonal projection onto \(E_j\),
and let \(B_{j,\varepsilon}\) be its second fundamental form.
These orthogonal projections need not be projections along the
specified complementary summands.

Write \(F_{V}\) for the Chern curvature of a Hermitian bundle \(V\).
The subbundle curvature formula gives
\[
 \tr\!\left(\sqrt{-1}\Lambda_{\omega_\varepsilon}F_{E_j}\right)
 =
 \tr\!\left(
   P_{j,\varepsilon}\sqrt{-1}\Lambda_{\omega_\varepsilon}F_{T_X}
 \right)-|B_{j,\varepsilon}|^2.
\]
Here and below the second fundamental forms have their complex
Hermitian norms.  This formula follows by writing the connection in
blocks for \(E_j\oplus E_j^\perp\); the off-diagonal blocks are
\(B_{j,\varepsilon}\) and \(-B_{j,\varepsilon}^*\).
In the Kähler case the mean curvature endomorphism
\(\sqrt{-1}\Lambda_{\omega_\varepsilon}F_{T_X}\) is the Ricci
endomorphism
\[
 H_\varepsilon=-I+\varepsilon A_\varepsilon,
 \qquad A_\varepsilon=\omega_\varepsilon^{-1}\omega_A>0.
\]

Integrating the trace formula computes first Chern degrees.
Since \(c_1(E_1)+c_1(E_2)=c_1(T_X)\), and writing
\(dV_\varepsilon=\omega_\varepsilon^n/n!\), we obtain
\[
 \int_X\sum_{j=1}^2|B_{j,\varepsilon}|^2\,dV_\varepsilon
 =
 \int_X\tr(C_\varepsilon H_\varepsilon)\,dV_\varepsilon,
 \qquad
 C_\varepsilon=P_{1,\varepsilon}+P_{2,\varepsilon}-I.
\]
The ranks add to \(n\), so \(\tr C_\varepsilon=0\).
Moreover \(-I\leq C_\varepsilon\leq I\), since each
\(P_{j,\varepsilon}\) is an orthogonal projection.  Positivity of
\(A_\varepsilon\) therefore gives
\begin{align}
 0
 &\leq\int_X\sum_j|B_{j,\varepsilon}|^2\,dV_\varepsilon
 \leq\varepsilon\int_X\tr(A_\varepsilon)\,dV_\varepsilon
 \notag\\
 &=\frac{\varepsilon}{(n-1)!}
       \int_X\omega_A\wedge\omega_\varepsilon^{n-1}
   =O(\varepsilon).
 \label{canon:energy}
\end{align}
The last integral is a bounded polynomial in \(\varepsilon\), by
Lemma~\ref{canon:metric-input}.  No angle bound between the summands
is used.

On compact subsets of \(X^0\), the metrics, their connections, the
orthogonal projections, and the second fundamental forms converge
smoothly.  A nonzero limiting second fundamental form would have a
fixed positive integral on some compact neighborhood, contradicting
\eqref{canon:energy}.  Both limiting second fundamental forms are
therefore zero.

The connection of a Kähler metric is torsion-free.  Hence each
summand is integrable on \(X^0\).  Its Frobenius obstruction
\[
 \Lambda^2E_j\longrightarrow T_X/E_j,\qquad
 u\wedge v\longmapsto [u,v]\bmod E_j,
\]
is a holomorphic bundle morphism on \(X\).  It vanishes on the
dense open set \(X^0\), hence everywhere.
Local Frobenius submersions for the two complementary foliations,
taken together, give holomorphic coordinates
\((z,w)\) in which the summands are the \(z\)- and \(w\)-directions.
Transitions between such charts locally preserve the two factors.
\end{proof}


\begin{proof}[Proof of Corollary~\ref{thm:nef-big}]
A smooth projective complex manifold is compact K\"ahler. By
Proposition~\ref{canon:parallel}, the supplied summands are integrable.
Theorem~\ref{thm:main} therefore applies and gives precisely the compatible
product claimed in the corollary.
\end{proof}

The argument just given uses the canonical metric only to establish
integrability. The appendices retain a second route to the global
product: the metric completion controls the size of crossings near an
entire exceptional fibre, while rational curves provide their holomorphic
continuation. The parallelism proved here is needed for that metric
argument, even though integrability alone also follows from
\cite[Theorem~1.3]{Horing07}.
